\chapter{DEPLOYMENT}

\section{Project Folder Structure}

The Vehicle Bidding System follows a standard full-stack application structure:

\begin{verbatim}
vehicle-bidding-system/
+-- frontend/
    +-- src/
        +-- components/          # Reusable React components
        +-- pages/               # Page components
        +-- context/             # React Context (AuthContext)
        +-- services/            # API service calls
        +-- App.jsx              # Main app component
        +-- main.jsx             # React entry point
        +-- index.css            # Global styles
    +-- public/                  # Static assets
    +-- package.json             # Frontend dependencies
    +-- vite.config.js           # Vite configuration
    +-- index.html               # HTML template
+-- backend/
    +-- server.js                # Express server entry point
    +-- middleware/
        +-- authMiddleware.js    # JWT authentication
    +-- models/                  # Mongoose schemas
        +-- User.js
        +-- Vehicle.js
        +-- Bid.js
        +-- Message.js
    +-- routes/                  # API route definitions
        +-- authRoutes.js
        +-- vehicleRoutes.js
        +-- bidRoutes.js
        +-- messageRoutes.js
        +-- reportRoutes.js
    +-- uploads/                 # Upload directory
    +-- package.json             # Backend dependencies
    +-- seedCars.js              # Sample data seeding
+-- .gitignore                   # Git ignore rules
+-- .env.example                 # Environment template
+-- README.md                    # Project documentation
\end{verbatim}

\section{Environment Setup}

\subsection{Prerequisites}

\begin{itemize}
    \item Node.js (v14.0 or higher)
    \item npm or yarn package manager
    \item MongoDB (local or MongoDB Atlas account)
    \item Git for version control
\end{itemize}

\subsection{Environment Variables}

Create a \texttt{.env} file in the backend directory with the following variables:

\begin{verbatim}
# Server Configuration
PORT=5000
NODE_ENV=development

# Database Configuration
MONGODB_URI=mongodb://localhost:27017/vehicle-bidding

# JWT Configuration
JWT_SECRET=your_jwt_secret_key_here
JWT_EXPIRY=7d

# Email Configuration (for password reset)
EMAIL_SERVICE=gmail
EMAIL_USER=your_email@gmail.com
EMAIL_PASSWORD=your_app_password

# File Upload Configuration
UPLOAD_DIR=./uploads
MAX_FILE_SIZE=5242880

# Frontend URL (for CORS)
FRONTEND_URL=http://localhost:5173

# Socket.io Configuration
SOCKET_CORS_ORIGIN=http://localhost:5173
\end{verbatim}

\section{Local Development Setup}

\subsection{Backend Setup}

\begin{itemize}
    \item Navigate to backend directory: \texttt{cd backend}
    \item Install dependencies: \texttt{npm install}
    \item Create \texttt{.env} file with configuration variables
    \item Start MongoDB service locally
    \item Run development server: \texttt{npm start} or \texttt{npm run dev}
    \item Server runs on \texttt{http://localhost:5000}
\end{itemize}

\subsection{Frontend Setup}

\begin{itemize}
    \item Navigate to frontend directory: \texttt{cd frontend}
    \item Install dependencies: \texttt{npm install}
    \item Create \texttt{.env.local} file if needed
    \item Start development server: \texttt{npm run dev}
    \item Application runs on \texttt{http://localhost:5173}
\end{itemize}

\section{Production Deployment}

\subsection{Docker Containerization}

Docker ensures consistency between development and production environments.

\subsubsection{Backend Docker Image}

Create a \texttt{Dockerfile} in the backend:

\begin{verbatim}
FROM node:18-alpine

WORKDIR /app

COPY package*.json ./
RUN npm install

COPY . .

EXPOSE 5000

CMD ["npm", "start"]
\end{verbatim}

\subsubsection{Frontend Docker Image}

Create a \texttt{Dockerfile} in the frontend:

\begin{verbatim}
FROM node:18-alpine as builder
WORKDIR /app
COPY package*.json ./
RUN npm install
COPY . .
RUN npm run build

FROM nginx:alpine
COPY --from=builder /app/dist /usr/share/nginx/html
COPY nginx.conf /etc/nginx/conf.d/default.conf
EXPOSE 80
CMD ["nginx", "-g", "daemon off;"]
\end{verbatim}

\subsection{Docker Compose}

Create a \texttt{docker-compose.yml} for orchestrating services:

\begin{verbatim}
version: '3.8'
services:
  mongodb:
    image: mongo:latest
    ports:
      - "27017:27017"
    environment:
      MONGO_INITDB_DATABASE: vehicle-bidding
    volumes:
      - mongo_data:/data/db

  backend:
    build: ./backend
    ports:
      - "5000:5000"
    environment:
      MONGODB_URI: mongodb://mongodb:27017/vehicle-bidding
    depends_on:
      - mongodb
    volumes:
      - ./backend:/app

  frontend:
    build: ./frontend
    ports:
      - "80:80"
    depends_on:
      - backend

volumes:
  mongo_data:
\end{verbatim}

\subsection{Cloud Deployment}

\subsubsection{Backend Deployment Options}

\begin{itemize}
    \item \textbf{Heroku:} Easy deployment with Git integration
    \item \textbf{AWS EC2:} Full control with scalability options
    \item \textbf{DigitalOcean:} Affordable and developer-friendly
    \item \textbf{Railway:} Simple deployment with Docker support
\end{itemize}

\subsubsection{Database Deployment}

\begin{itemize}
    \item \textbf{MongoDB Atlas:} Managed MongoDB in the cloud
    \item \textbf{AWS DocumentDB:} MongoDB-compatible database
    \item \textbf{Self-hosted:} MongoDB on dedicated server
\end{itemize}

\subsubsection{Frontend Deployment}

\begin{itemize}
    \item \textbf{Vercel:} Optimized for React applications
    \item \textbf{Netlify:} Simple deployment with CI/CD
    \item \textbf{AWS S3 + CloudFront:} Cost-effective with CDN
    \item \textbf{GitHub Pages:} Free for static sites
\end{itemize}

\subsection{Deployment Checklist}

Before deploying to production:

\begin{itemize}
    \item Set \texttt{NODE\_ENV=production}
    \item Configure all environment variables
    \item Enable HTTPS/SSL certificates
    \item Configure CORS properly
    \item Set up database backups
    \item Configure logging and monitoring
    \item Test all API endpoints
    \item Verify email functionality
    \item Test payment integration (if applicable)
    \item Run security audit
    \item Configure load balancing if needed
\end{itemize}

\section{Build and Packaging}

\subsection{Frontend Build}

\begin{verbatim}
npm run build
\end{verbatim}

This creates an optimized production build in the \texttt{dist} directory.

\subsection{Backend Build}

The backend doesn't require building; it runs directly with Node.js. However, you can use bundlers like:

\begin{itemize}
    \item Webpack for bundle optimization
    \item Rollup for library-style builds
    \item ESBuild for fast builds
\end{itemize}

\section{Performance Optimization}

\subsection{Frontend Optimization}

\begin{itemize}
    \item Code splitting with dynamic imports
    \item Image optimization and lazy loading
    \item CSS minification and critical path CSS
    \item JavaScript minification and tree shaking
    \item Gzip compression
    \item Browser caching strategies
\end{itemize}

\subsection{Backend Optimization}

\begin{itemize}
    \item Database indexing on frequently queried fields
    \item Query optimization
    \item Caching with Redis
    \item Compression middleware
    \item Connection pooling
    \item Load balancing across servers
\end{itemize}

\section{Monitoring and Maintenance}

\subsection{Logging}

Implement logging at multiple levels:

\begin{itemize}
    \item Application logs (errors, warnings, info)
    \item API request/response logs
    \item Database query logs
    \item Performance metrics
\end{itemize}

\subsection{Monitoring Tools}

\begin{itemize}
    \item \textbf{PM2:} Node.js process management
    \item \textbf{NewRelic:} Application performance monitoring
    \item \textbf{Sentry:} Error tracking
    \item \textbf{Datadog:} Infrastructure monitoring
    \item \textbf{CloudWatch:} AWS monitoring service
\end{itemize}

\subsection{Backup and Recovery}

\begin{itemize}
    \item Daily automated database backups
    \item Backup retention policy (30 days minimum)
    \item Regular backup restoration testing
    \item Disaster recovery plan
    \item Data redundancy and replication
\end{itemize}

\section{Security Deployment Considerations}

\begin{itemize}
    \item Use HTTPS/TLS with valid SSL certificates
    \item Implement CORS headers properly
    \item Use security headers (HSTS, CSP, X-Frame-Options)
    \item Regular security updates and patches
    \item Implement rate limiting and DDoS protection
    \item Regular security audits and penetration testing
    \item Secure credential management
    \item Database encryption (at rest and in transit)
    \item Regular backups with encryption
\end{itemize}

\section{Scaling Strategies}

\subsection{Horizontal Scaling}

\begin{itemize}
    \item Load balancer (Nginx, HAProxy)
    \item Multiple backend server instances
    \item Session management across servers
    \item Database replication and sharding
\end{itemize}

\subsection{Vertical Scaling}

\begin{itemize}
    \item Upgrading server resources (CPU, RAM)
    \item Optimization of existing resources
    \item Connection pooling and caching
\end{itemize}

\section{Troubleshooting}

\subsection{Common Issues}

\begin{itemize}
    \item \textbf{Database Connection:} Verify MongoDB URI and authentication
    \item \textbf{CORS Errors:} Check frontend URL in backend configuration
    \item \textbf{Socket.io Connection:} Verify Socket.io configuration and CORS
    \item \textbf{File Upload Issues:} Check upload directory permissions
    \item \textbf{JWT Token Expiry:} Verify token refresh mechanism
\end{itemize}

\subsection{Support and Maintenance}

\begin{itemize}
    \item Regular code reviews
    \item Security patch management
    \item Dependency updates and upgrades
    \item Documentation updates
    \item Team knowledge sharing sessions
\end{itemize}

